\documentclass[a4paper,11pt]{article}
\usepackage[utf8]{inputenc}
\usepackage[french]{babel}
\usepackage[T1]{fontenc}
\usepackage{amsmath,amssymb,amsthm}
\usepackage{geometry}
\geometry{margin=2cm}
\usepackage{enumitem}
\usepackage{hyperref}
\usepackage{booktabs}
\usepackage{array}
\usepackage{xcolor}
\usepackage{listings}
\usepackage{titlesec}
\usepackage{fancyhdr}
\usepackage{lastpage}
\usepackage{graphicx}
\usepackage{etoolbox}
\AtBeginDocument{\shorthandoff{;:!?}}

\titleformat{\section}{\large\bfseries\color{blue!60!black}}{}{0pt}{}[\vspace{-0.5ex}\rule{\textwidth}{0.4pt}]
\titleformat{\subsection}{\bfseries\color{blue!40!black}}{}{0pt}{}
\titleformat{\subsubsection}{\itshape\color{blue!30!black}}{}{0pt}{}

\lstdefinestyle{monstyle}{
  frame=single,
  numbers=left,
  numbersep=5pt,
  breaklines=true,
  basicstyle=\small\ttfamily,
  backgroundcolor=\color{gray!5},
  rulecolor=\color{gray!40},
  columns=fullflexible,
  keepspaces=true,
  showstringspaces=false
}
\lstset{style=monstyle}

\pagestyle{fancy}
\fancyhf{}
\fancyhead[L]{\small STA211 -- Isolation Forest}
\fancyhead[R]{\small Fiche r\'ecapitulative}
\fancyfoot[C]{\thepage/\pageref{LastPage}}
\renewcommand{\headrulewidth}{0.4pt}

\theoremstyle{definition}
\newtheorem*{definition}{D\'efinition}
\theoremstyle{remark}
\newtheorem*{remark}{Remarque}
\newtheorem*{important}{\`A retenir}

\title{\textbf{Fiche r\'ecapitulative -- STA211}\\[4pt]
\large D\'etection des points atypiques via \emph{Isolation Forest}}
\author{}
\date{10 juin 2026}

\begin{document}
\maketitle
\thispagestyle{empty}

\begin{abstract}
Cette fiche pr\'esente la m\'ethode \emph{Isolation Forest} pour la d\'etection
non supervis\'ee des valeurs atypiques. Elle couvre le principe (isolement
rapide des anomalies par des arbres al\'eatoires), la construction des arbres,
le calcul du score d'anomalie, le param\'etrage, et une application sur les
donn\'ees German Credit avec le package R \texttt{isotree}.
\end{abstract}

\vfill
\tableofcontents
\newpage

% ===================================================================
\section{Contexte}
% ===================================================================

La d\'etection des points atypiques (ou \emph{outliers}) est une \'etape
importante en fouille de donn\'ees. Du point de vue g\'eom\'etrique, il s'agit
de points qui s'\'ecartent significativement des autres~: ils sont peu nombreux
et localis\'es dans une zone peu dense du nuage de points. Leur d\'etection est
un enjeu car les analyses effectu\'ees sur des \'echantillons contenant ces
points peuvent conduire \`a des r\'esultats tr\`es diff\'erents.

\emph{Isolation Forest} (\cite{LiuTZ08}) est une m\'ethode de r\'ef\'erence
qui s'appuie sur les for\^ets al\'eatoires dans un cadre non supervis\'e.

% ===================================================================
\section{Principe}
% ===================================================================

La d\'etection se fait ici d'un point de vue compl\`etement
\textbf{non supervis\'e}. Les for\^ets \'etant des m\'ethodes supervis\'ees,
la construction des arbres y est diff\'erente.

Le principe d'\emph{Isolation Forest}~:
\begin{enumerate}
  \item Construire des arbres \textbf{au hasard} qui vont fournir
    diff\'erents d\'ecoupages de l'espace.
  \item Le postulat~: les valeurs atypiques vont rapidement \^etre isol\'ees
    dans une r\'egion (peu de coupures n\'ecessaires), contrairement aux
    observations normales qui seront isol\'ees \`a un niveau plus profond
    de l'arbre.
  \item On regarde \`a quelle profondeur de l'arbre une observation est
    isol\'ee en moyenne. Ceci fournit un score par observation.
  \item Les scores sont ensuite seuill\'es pour d\'eterminer les valeurs
    atypiques.
\end{enumerate}

\begin{important}
  L'id\'ee centrale~: plus une observation est isol\'ee rapidement (faible
  profondeur), plus elle est susceptible d'\^etre atypique.
\end{important}

% ===================================================================
\section{Construction des arbres au hasard}
% ===================================================================

\subsection{Algorithme de construction d'un arbre}

\begin{enumerate}
  \item S\'electionner au hasard une variable parmi les $p$ variables.
  \item Choisir au hasard un point de coupure dans l'\'etendue de cette
    variable.
  \item Segmenter l'\'echantillon en deux sous-groupes.
  \item R\'ep\'eter it\'erativement jusqu'\`a ce qu'un n\oe ud~:
  \begin{enumerate}
    \item ne contienne qu'un seul individu, ou
    \item atteigne une profondeur pr\'ed\'efinie.
  \end{enumerate}
\end{enumerate}

\subsection{Sous-\'echantillonnage et effet de masque}

Pour construire une for\^et, la proc\'edure commence par un
\textbf{sous-\'echantillonnage} afin de limiter \textbf{l'effet de masque}
(\emph{masking effect}). Ce ph\'enom\`ene correspond au fait que la
proc\'edure de d\'etection doit elle-m\^eme faire face aux valeurs atypiques.
Les valeurs atypiques peuvent avoir un impact sur la m\'ethode de d\'etection
et ainsi masquer certaines d'entre elles.

En tirant un sous-\'echantillon, les points atypiques seront plus rapidement
isol\'es et donc mieux d\'etect\'es.

\subsection{For\^et}

Une fois le sous-\'echantillonnage effectu\'e, on construit un premier arbre
au hasard, puis on recommence $B$ fois pour obtenir une for\^et de $B$ arbres.

% ===================================================================
\section{Calcul du score}
% ===================================================================

\subsection{Profondeur d'isolement}

\'Etant donn\'ee une observation $x$, on note $h(x)$ la profondeur de
$x$ dans un arbre (nombre d'ar\^etes entre la racine et le n\oe ud
terminal contenant $x$).

Sur la base d'une for\^et de $B$ arbres, on note
$\overline{h}(x) = \frac{1}{B} \sum_{b=1}^{B} h_b(x)$ la profondeur
moyenne de $x$ dans la for\^et.

\subsection{Normalisation par $c(n)$}

Comme la profondeur d'un arbre d\'epend du nombre d'observations, on
relativise $\overline{h}(x)$ par la profondeur moyenne attendue pour un
arbre de $n$ observations~:

\begin{equation}
c(n) = 2 \times H(n-1) - \frac{2(n-1)}{n}
\end{equation}

o\`u $H(k) = \sum_{i=1}^{k} \frac{1}{i}$ est le $k$-i\`eme nombre
harmonique. $c(n)$ est la profondeur moyenne d'un arbre construit au
hasard sur $n$ observations.

\subsection{Score d'anomalie}

On d\'efinit le score de l'observation $x$~:

\begin{equation}
s(x, n) = 2^{-\frac{\overline{h}(x)}{c(n)}}
\end{equation}

\textbf{Interpr\'etation~:}
\begin{itemize}
  \item $s(x, n) \approx 0.5$~: l'observation est normale
    ($\overline{h}(x) \approx c(n)$).
  \item $s(x, n) \to 1$~: l'observation est tr\`es probablement atypique
    ($\overline{h}(x) \ll c(n)$, isolement rapide).
  \item $s(x, n) \to 0$~: l'observation est plut\^ot normale
    ($\overline{h}(x) \gg c(n)$, isolement profond).
\end{itemize}

\begin{important}
  Le score $s(x,n)$ est une standardisation de $\overline{h}(x)$ comprise
  entre 0 et 1. Plus le score est proche de 1, plus l'observation est
  atypique.
\end{important}

% ===================================================================
\section{Param\'etrage}
% ===================================================================

La m\'ethode d\'epend de trois param\`etres principaux~:

\subsection{Nombre d'arbres ($B$)}

Le nombre d'arbres a une incidence sur la variance des profondeurs
moyennes $\overline{h}(x)$. Plus $B$ est grand, plus l'estimation est
stable. En pratique, $B = 100$ est souvent suffisant, $B = 300$ donne
une tr\`es bonne stabilit\'e.

\subsection{Taille du sous-\'echantillon}

\cite{LiuTZ08} indiquent que tirer un sous-\'echantillon de
\textbf{256 observations} est suffisant et qu'il n'y a pas de gain
substantiel \`a choisir un nombre plus \'elev\'e. Cela permet de r\'eduire
le co\^ut computationnel.

\begin{remark}
  Ce nombre de 256 est souvent repris dans la litt\'erature mais sa
  port\'e reste limit\'ee (bas\'e sur une \'etude par simulation).
\end{remark}

\subsection{Seuillage des scores}

Une fois le score calcul\'e pour chaque observation, il convient de le
seuiller pour d\'eterminer les observations atypiques. Il n'y a pas de
valeur seuil universelle.

M\'ethode pratique~: observer la distribution des scores et chercher
un d\'ecrochage (\emph{elbow}) dans la distribution. Les points au-del\`a
du coude sont consid\'er\'es comme atypiques.

% ===================================================================
\section{Application aux donn\'ees German Credit}
% ===================================================================

On utilise le package \texttt{isotree} (\cite{Cortes2026}) sur les donn\'ees
German Credit ($n=1000$, $p=21$, pas de valeurs manquantes).

\subsection{Chargement et importation}

\begin{lstlisting}[language=R]
library(isotree)
library(FactoMineR)

don <- read.table(
  "https://archive.ics.uci.edu/ml/machine-learning-databases
   /statlog/german/german.data",
  sep = " ", stringsAsFactors = TRUE)
\end{lstlisting}

\subsection{Construction du mod\`ele}

On applique \emph{Isolation Forest} avec $B = 300$ arbres et un
sous-\'echantillon de taille 256~:

\begin{lstlisting}[language=R]
set.seed(1234)
ntree <- 300
sample_size <- 256

model <- isolation.forest(don,
  sample_size = sample_size,
  ntrees = ntree)
\end{lstlisting}

\subsection{Stabilisation des scores}

On \'etudie la stabilisation des scores en fonction du nombre d'arbres~:

\begin{lstlisting}[language=R]
subscores <- sapply(seq.int(ntree),
  FUN = function(nbtree, model, don) {
    submodel <- isotree.subset.trees(model,
      trees_take = seq.int(nbtree))
    predict.isolation_forest(submodel, newdata = don)
  }, model = model, don = don)

matplot(t(subscores), type = "b", pch = 14, cex = .1,
  ylab = "score", xlab = "nombre d'arbres")
\end{lstlisting}

Les scores se stabilisent \`a partir de 100 arbres.

\subsection{Seuillage et identification}

\begin{lstlisting}[language=R]
scores <- predict.isolation_forest(model, newdata = don)
scores <- sort(scores, decreasing = TRUE)

plot(scores, type = "b", pch = 16, cex = 0.2,
     xlab = "rang")
threshold <- 0.576
abline(h = threshold, col = "red", lty = 2)

which(scores > threshold)
# 375 757 916 187 851 227 928 890 614
\end{lstlisting}

Neuf observations sont identifi\'ees comme atypiques avec un seuil
de 0.576.

\subsection{Post-analyse}

L'\emph{Isolation Forest} fournit un score par observation sur la base
de l'ensemble des variables, mais n'indique pas quelles variables
expliquent ce caract\`ere atypique. On peut utiliser une analyse
factorielle (AFDM) pour analyser les caract\'eristiques des individus
atypiques~:

\begin{lstlisting}[language=R]
res.FAMD <- FAMD(don,
  ind.sup = as.numeric(names(which(scores > threshold))),
  sup.var = which(colnames(don) == "Creditability"),
  graph = FALSE)
plot(res.FAMD, choix = "ind", cex = .5)
\end{lstlisting}

Il s'agit ensuite d'identifier les variables contribuant \`a la
construction des axes sur lesquels ces observations ont des
coordonn\'ees \'elev\'ees.

% ===================================================================
\section{Conclusion}
% ===================================================================

\emph{Isolation Forest} est une m\'ethode de r\'ef\'erence pour la
d\'etection des valeurs atypiques~:

\begin{itemize}
  \item Faible co\^ut algorithmique (complexit\'e lin\'eaire),
    permettant de traiter des \'echantillons de grande taille.
  \item Non sensible aux probl\`emes d'\'echelle (pas de
    standardisation n\'ecessaire).
  \item Extension aux donn\'ees mixtes directe.
\end{itemize}

Une fois les points atypiques identifi\'es, on peut~:
\begin{itemize}
  \item Utiliser des m\'ethodes d'analyse \emph{robustes}.
  \item Remplacer les valeurs jug\'ees atypiques par des valeurs
    manquantes puis effectuer une imputation.
\end{itemize}

Le package \texttt{isotree} (\cite{Cortes2026}) propose diff\'erentes
extensions, notamment la segmentation selon des combinaisons de
variables.

% ===================================================================
\section*{Questions cl\'es (QCM)}
% ===================================================================

\addcontentsline{toc}{section}{Questions cl\'es (QCM)}

\begin{enumerate}

  \item Quel est le principe de base d'\emph{Isolation Forest}~?
    \begin{enumerate}
      \item Les points atypiques sont isol\'es rapidement par les arbres
      \item Les points atypiques sont identifi\'es par leur distance au centre
      \item Les points atypiques sont d\'etect\'es par des m\'ethodes
        supervis\'ees
      \item Les points atypiques sont toujours en profondeur maximale
    \end{enumerate}
    \textbf{R\'eponse~: a (isolement rapide des anomalies).}

  \item Dans la construction d'un arbre, comment sont choisis la variable
    et le point de coupure~?
    \begin{enumerate}
      \item Par optimisation d'un crit\`ere (Gini, entropie)
      \item Au hasard
      \item Par validation crois\'ee
      \item Par recherche exhaustive
    \end{enumerate}
    \textbf{R\'eponse~: b (variable et point de coupure choisis au hasard).}

  \item Quelle est la formule du score d'anomalie $s(x,n)$~?
    \begin{enumerate}
      \item $s(x,n) = 2^{-\overline{h}(x)/c(n)}$
      \item $s(x,n) = \overline{h}(x)/c(n)$
      \item $s(x,n) = 2^{\overline{h}(x)/c(n)}$
      \item $s(x,n) = -\overline{h}(x)/c(n)$
    \end{enumerate}
    \textbf{R\'eponse~: a ($s(x,n) = 2^{-\overline{h}(x)/c(n)}$).}

  \item Que repr\'esente $c(n)$ dans le calcul du score~?
    \begin{enumerate}
      \item La profondeur maximale de l'arbre
      \item La profondeur moyenne attendue pour un arbre al\'eatoire
        sur $n$ observations
      \item Le nombre d'arbres dans la for\^et
      \item La variance des scores
    \end{enumerate}
    \textbf{R\'eponse~: b (profondeur moyenne attendue).}

  \item Un score $s(x,n) \approx 0.5$ indique~:
    \begin{enumerate}
      \item Une observation tr\`es atypique
      \item Une observation normale
      \item Une observation anormale
      \item Une erreur de calcul
    \end{enumerate}
    \textbf{R\'eponse~: b ($\overline{h}(x) \approx c(n)$, observation
    normale).}

  \item Quelle est la taille de sous-\'echantillon recommand\'ee par
    Liu, Ting \& Zhou (2008)~?
    \begin{enumerate}
      \item 128
      \item 256
      \item 512
      \item 1000
    \end{enumerate}
    \textbf{R\'eponse~: b (256 observations).}

  \item Quel est l'int\'er\^et du sous-\'echantillonnage dans
    \emph{Isolation Forest}~?
    \begin{enumerate}
      \item R\'eduire la variance du mod\`ele
      \item Limiter l'effet de masque
      \item Acc\'el\'erer l'entra\^inement seulement
      \item \'Eviter le sur-apprentissage
    \end{enumerate}
    \textbf{R\'eponse~: b (limiter l'effet de masque).}

  \item Quel package R impl\'emente \emph{Isolation Forest}~?
    \begin{enumerate}
      \item randomForest
      \item isotree
      \item rpart
      \item xgboost
    \end{enumerate}
    \textbf{R\'eponse~: b (\texttt{isotree}).}

  \item \emph{Isolation Forest} n\'ecessite-t-elle une standardisation
    pr\'ealable des donn\'ees~?
    \begin{enumerate}
      \item Oui, toujours
      \item Non, elle n'est pas sensible aux \'echelles
      \item Oui, si les variables sont de natures diff\'erentes
      \item Cela d\'epend du jeu de donn\'ees
    \end{enumerate}
    \textbf{R\'eponse~: b (non sensible aux \'echelles, pas de
    standardisation n\'ecessaire).}

  \item Dans l'application German Credit, combien d'observations sont
    identifi\'ees comme atypiques avec le seuil de 0.576~?
    \begin{enumerate}
      \item 3
      \item 6
      \item 9
      \item 12
    \end{enumerate}
    \textbf{R\'eponse~: c (9 observations atypiques).}

\end{enumerate}

% ===================================================================
\begin{thebibliography}{9}
% ===================================================================

\bibitem{LiuTZ08}
  F.~T.~Liu, K.~M.~Ting, Z.-H.~Zhou.
  \newblock \emph{Isolation Forest}.
  \newblock ICDM, pp.~413--422, IEEE, 2008.

\bibitem{Cortes2026}
  D.~Cortes.
  \newblock \emph{isotree: Isolation-Based Outlier Detection}.
  \newblock \url{https://doi.org/10.32614/CRAN.package.isotree}, 2026.

\bibitem{Audigier2026}
  V.~Audigier.
  \newblock STA211~: d\'etection des points atypiques via
    \emph{Isolation Forest}.
  \newblock Cours CNAM, 2026.

\bibitem{Rakotomalala2026}
  R.~Rakotomalala.
  \newblock D\'etection des anomalies~: \emph{Isolation Forest}.
  \newblock \url{https://eric.univ-lyon2.fr/ricco/cours/slides/isolation_forest.pdf}

\end{thebibliography}

\end{document}
